% Part: first-order-logic
% Chapter: syntax-and-semantics
% Section: subformulas

\documentclass[../../../include/open-logic-section]{subfiles}

\begin{document}

\olfileid{fol}{syn}{sbf}

\olsection{\printtoken{P}{subformula}}

\begin{explain}
كثيرًا ما يكون من المفيد أن نتحدث عن ال!!{formula}s التي «تؤلف»
!!a{formula} معطاة. ونسميها \emph{!!{subformula}s} لها. وتُعَد كل
!!{formula} !!a{subformula} من نفسها؛ أما الصيغة الجزئية من~$!A$
المغايرة لـ~$!A$ نفسها فهي \emph{!!{subformula} حقيقية}.
\end{explain}

\begin{defn}[!!^{subformula} مباشرة]
إذا كانت $!A$ !!a{formula}، فتُعرّف \emph{ال!!{subformula}s
المباشرة} لـ~$!A$ استقرائيًا كما يأتي:
\begin{enumerate}
\item ليس لل!!{formula}s الذرية !!{subformula}s مباشرة.

\tagitem{prvNot}{\indcase{!A}{\lnot !B}{ال!!{subformula} المباشرة
    الوحيدة لـ~$\indfrm$ هي~$!B$.}}{}

\item \indcase{!A}{(!B \ast !C)}{ال!!{subformula}s المباشرة
  لـ~$\indfrm$ هي $!B$ و$!C$ (حيث $\ast$ أي رابط ثنائي).}

\tagitem{prvAll}{\indcase{!A}{\lforall[x][!B]}{ال!!{subformula}
    المباشرة الوحيدة لـ~$\indfrm$ هي~$!B$.}}{}

\tagitem{prvEx}{\indcase{!A}{\lexists[x][!B]}{ال!!{subformula}
    المباشرة الوحيدة لـ~$\indfrm$ هي~$!B$.}}{}
\end{enumerate}
\end{defn}

\begin{defn}[!!^{subformula} حقيقية]
إذا كانت $!A$ !!a{formula}، فتُعرّف \emph{ال!!{subformula}s
الحقيقية} لـ~$!A$ عوديًا كما يأتي:
\begin{enumerate}
\item ليس لل!!{formula}s الذرية !!{subformula}s حقيقية.

\tagitem{prvNot}{\indcase{!A}{\lnot !B}{ال!!{subformula}s الحقيقية
    لـ~$\indfrm$ هي~$!B$ مع جميع ال!!{subformula}s الحقيقية
    لـ~$!B$.}}{}

\item \indcase{!A}{(!B \ast !C)}{ال!!{subformula}s الحقيقية
  لـ~$\indfrm$ هي $!B$ و$!C$، مع جميع ال!!{subformula}s الحقيقية
  لـ~$!B$ وتلك الخاصة بـ~$!C$.}

\tagitem{prvAll}{\indcase{!A}{\lforall[x][!B]}{ال!!{subformula}s
    الحقيقية لـ~$\indfrm$ هي~$!B$ مع جميع ال!!{subformula}s الحقيقية
    لـ~$!B$.}}{}

\tagitem{prvEx}{\indcase{!A}{\lexists[x][!B]}{ال!!{subformula}s
    الحقيقية لـ~$\indfrm$ هي~$!B$ مع جميع ال!!{subformula}s الحقيقية
    لـ~$!B$.}}{}
\end{enumerate}
\end{defn}

\begin{defn}[!!^{subformula}]
ال!!{subformula}s لـ~$!A$ هي $!A$ نفسها مع جميع ال!!{subformula}s
الحقيقية لها.
\end{defn}

\begin{explain}
لاحظ الفرق الدقيق بين طريقتنا في تعريف ال!!{subformula}s المباشرة
وطريقتنا في تعريف ال!!{subformula}s الحقيقية. ففي الحالة الأولى
عرّفنا مباشرة ال!!{subformula}s المباشرة لصيغة~$!A$ في كل صورة ممكنة
لـ~$!A$. وهذا تعريف صريح بالحالات، وحالاته تحاكي التعريف الاستقرائي
لمجموعة ال!!{formula}s. وفي الحالة الثانية حاكينا أيضًا طريقة تعريف
مجموعة جميع ال!!{formula}s، لكننا أدرجنا في كل حالة ال!!{subformula}s
الحقيقية لل!!{formula}s الأصغر $!B$ و$!C$، فضلًا عن هاتين
ال!!{formula}s نفسيهما. وهذا يجعل التعريف \emph{عوديًا}. وبوجه عام
يكون تعريف دالة على مجموعة معرّفة استقرائيًا (وهي في حالتنا
ال!!{formula}s) عوديًا إذا استخدمت حالات تعريف الدالة الدالة نفسها.
ولكن لكي يكون التعريف سليمًا، يجب أن نتأكد من أننا لا نستخدم قط إلا
قيم الدالة عند وسائط «تسبق» الوسيط الذي نعرّفه---ففي حالتنا، حين
نعرّف «!!{subformula} حقيقية» لـ~$(!B \ast !C)$، لا نستخدم إلا
ال!!{subformula}s الحقيقية لل!!{formula}s «الأسبق» $!B$ و$!C$.
\end{explain}

\begin{prop}
\ollabel{prop:subfrm-trans}
افترض أن $!B$ صيغة جزئية من~$!A$، وأن $!C$ صيغة جزئية من~$!B$.
عندئذ تكون $!C$ صيغة جزئية من~$!A$. وبعبارة أخرى، علاقة الصيغة
الجزئية متعدية.
\end{prop}

\begin{prob}
أثبت \olref[fol][syn][sbf]{prop:subfrm-trans}.
\end{prob}

\begin{prop}
\ollabel{prop:count-subfrms}
افترض أن $!A$ صيغة فيها $n$ من الروابط والمكممات. عندئذ يكون لـ~$!A$
على الأكثر $2n+1$ من الصيغ الجزئية.
\end{prop}

\begin{prob}
أثبت \olref[fol][syn][sbf]{prop:count-subfrms}.
\end{prob}

\end{document}
